\section{Reproducing the finite results}\label{app:reproduction}

The artifact provides one command that checks the direct construction's
headline results. From the repository root, run
\begin{quote}
\small\ttfamily
python field-universality/paper/reproduce\_results.py
\end{quote}
The companion file \texttt{paper/REPRODUCING.md}, under
\texttt{field-universality/}, gives dependencies, certificate paths,
build commands, and the procedure for repeating the timing runs.
Python assertions must remain enabled: the verifiers rely on
\texttt{assert} statements, so the runner refuses to start under
\texttt{python -O}. The runner reads a fixed list of certificates, the
publication selection, recorded in
\path{results/distance-time/certificate-audit.json} under
\path{field-universality/expander-lift/}; it never searches for newer
certificates. The inventory \path{paper/results-inventory.json} under
\path{field-universality/}, generated by
\path{paper/build_result_inventory.py}, records for every
reported point its distance, failure bound, timing, and the hashes of
its evidence files. Tables~\ref{tab:width-one-small},
\ref{tab:wide-measured}, and~\ref{tab:numerical-grid} are generated
from the same records.

\paragraph{What the runner checks.}
The runner performs four checks. First, it verifies twelve finite
certificates: the eleven measured Newton points in
Tables~\ref{tab:width-one-small} and~\ref{tab:wide-measured}, and the
subset-sum point of Theorem~\ref{thm:subset-sum-gv}. For each, it
recomputes the verified summary from the certificate and compares it
with the saved record. Second, it reconstructs the deterministic
ceilings of Section~\ref{sec:group-ceiling} by exact integer counting.
Third, it checks the archived benchmark source hashes and correctness
logs, recomputes the timing summaries, and checks each timing's
association with its certificate and with the plotting data. Fourth,
it regenerates the inventory and all three generated tables, including
the numerical grid, and compares them with the manuscript files. The
grid entries are compared with their stored records; they are not
reoptimized. These checks verify numerical bounds under the paper's
mathematical arguments; they are not a formal verification of those
arguments. Checking archived logs does not rerun the encoder or
reproduce the timing experiment on another host.

\paragraph{Certificates and verifiers.}
Unless stated otherwise, paths are relative to
\path{field-universality/expander-lift/}; certificate files without a
directory are in \path{results/}, and verifier scripts without a
directory are in the top level. Results marked $\dagger$ are checked by
the runner; the others have separate verifiers.
\begin{itemize}
\item Width-one Newton points, $q=1$ (Appendix~\ref{app:complete-runs})$^\dagger$:
  \path{all-one-newton-line-adaptive-e245366.json}, checked by
  \path{all_one_line_run_certificate.py}.
\item Width-one Newton points, $q=2,\ldots,256$
  (Appendix~\ref{app:small-q-pairs})$^\dagger$: the rows of
  \path{results/distance-time/certificate-audit.json}, checked by
  \path{small_q_fresh_pair_certificate.py}.
\item Widths two and four (Appendix~\ref{app:wide-mask})$^\dagger$:
  the rows of the same manifest under
  \path{results/convolution-wide-finite/}, checked by
  \path{convolution_wide_finite.py}.
\item Subset-sum precode, Theorem~\ref{thm:subset-sum-gv}$^\dagger$:
  \path{all-one-line-runs-e263402.json}, checked by
  \path{all_one_line_run_certificate.py}, including reconstruction of
  the exact local catalog.
\item Subset-sum precode with the joint-renewal bound,
  $e=239075$ (Appendix~\ref{app:subset-sum-gv}):
  \path{all-one-binary-hadamard-e239075.json}, checked by
  \path{all_one_hadamard_certificate.py}.
\item Deterministic ceilings (Section~\ref{sec:group-ceiling})$^\dagger$:
  \path{results/convolution-width-one-ceilings.json}, checked by
  \path{convolution_width_one_ceiling.py}.
\item Two-state bound, $e=138412$ (Appendix~\ref{app:all-one}):
  \path{results/convolution-width-one-direct-r8-q1-e138412.json},
  checked by \path{convolution_width_one_finite.py}.
\item Projective-pair bound, $e=209715$ (Appendix~\ref{app:fresh-pairs}):
  \path{results/all-one-projective-pairs-r8-e209715.json}, checked by
  \path{all_one_projective_pairs.py}.
\item Joint-renewal bound, $e=224395$ (Appendix~\ref{app:joint-renewal}):
  \path{results/all-one-joint-renewal-r8-e224395.json}, checked by
  \path{all_one_renewal_certificate.py}.
\item Four-input precode, Theorem~\ref{thm:small-scans}
  (Appendix~\ref{app:small-certificate}): the summary
  \path{results/scalar-frontier/summary.json} and the certificates
  \path{results/scalar-frontier/distance-q2048-e304087.json} and
  \path{results/scalar-frontier/margin-q2048-e230896.json}, generated by
  \path{scalar_frontier.py} and checked by \path{distance_ladder.py},
  which verifies the rational side conditions, interval endpoints,
  coverage, sparse sum, total failure bound, and a $2^{-40}$ budget.
  The tests in \path{test_scalar_frontier.py} check every complete
  certificate in the summary. The $q=32$ row is in
  \path{results/signed-scan-certificate.json}, with verifier
  \path{signed_scan_certificate.py}.
\item Four-input timings (Table~\ref{tab:small-timing}): the archive
  \path{results/scalar-alphabet-benchmark/}, with runner
  \path{bench/run_scalar_alphabet.sh} and archive checker
  \path{bench/summarize_scalar_alphabet.py}. The archive includes the
  measured sources, raw timing samples, host and compiler details,
  hashes, and validation outputs.
\item Fixed-prime BAA bounds, Theorem~\ref{thm:baa-fixed}
  (Appendix~\ref{app:fixed-prime}); paths relative to
  \path{field-universality/}: the summaries
  \path{results/ternary-bulk-0.196-eta-0.0001-independent.json},
  \path{results/quinary-sign-orbit-bulk-0.05-eta-1-over-60.json}, and
  \path{results/septic-sign-orbit-bulk-0.018-eta-0.0625.json},
  reconstructed and checked by
  \path{scripts/verify_symmetric_ternary.py} and
  \path{scripts/verify_sign_orbit_bulk.py}. Each summary's
  \texttt{command} field gives the full cover reconstruction command;
  reading a summary alone does not check the cover.
\item Walsh--Hadamard union bound~\eqref{eq:wht-three-flat} and the
  sparse BAA parameter rows: \path{paper/audit-supporting-checks.json},
  relative to \path{field-universality/}, written by
  \path{paper/check_audit_arithmetic.py} using rational arithmetic.
\end{itemize}

\paragraph{Measured curve.}
The timing summaries come from 72 processes and 2232 timed encodes:
twelve kernels, each run with six setup seeds and 31 timed encodes per
seed. For $q=2$ there are two implementations, one with additions and
one with small-integer multiplications; they produce identical output
checksums, and the figure and tables use the one with the smaller
median time. The reported latency for a kernel is the median of its
six process medians, that is, the mean of the third and fourth; each
process median is the median of 31 timed encodes.
Figure~\ref{fig:distance-time} shows the range of those process medians.

\paragraph{Measurement environment.}
All timings used one core of an AMD Ryzen 9 7950X, pinned to CPU 2.
For the Newton encoders, boost was disabled with a 4.5 GHz frequency
setting, and GCC 15.2.0 compiled with \texttt{-O3 -march=native
-DNDEBUG} and C++20. Each configuration used two warmups before its
timed encodes. Runs were serial, with case order reversed on alternate
seeds. Setup, allocation, reference checks, and output checksums were
outside the timer. Setup uses reproducible pseudorandom streams; the
distance theorems concern the specified ideal sampling distribution.
Complete outputs were checked against independent references and
integer arithmetic oracles, and sanitizer runs passed. The four-input
timings of Table~\ref{tab:small-timing} used the same processor and
core, GCC 15.2 with \texttt{-O3 -march=native}, 42 processes, and 1302
timed encodes. Its pages were resident and caches were not flushed.
Its outputs matched a dense-local generic-field reference; boundary and
tail tests, multiprecision arithmetic checks through coefficient 2048,
and address and undefined-behavior sanitizers passed. SplitMix64
generated its reproducible benchmark fixtures; these fixtures are not a
cryptographic setup implementation.

\paragraph{Scope of the comparison.}
All eleven measured points use the Newton $[16,8,9]$ precode $B_8$,
two permutations, no final unpermuted scan, and at least forty bits of
setup-failure margin. The subset-sum point uses a different precode and
a smaller margin, and was not timed. Width eight has numerical
dense-support estimates only; it is not part of the finite-certificate
or timing checks. The runner does not check the BAA and Walsh--Hadamard
case studies (Sections~\ref{sec:baa} and~\ref{sec:wht}), the separate
verifiers listed above, or any search for optimal parameters.
